\section{合同的保全}
\begin{enumerate}
  \item 必考。主客观都考，分多。担保制度的一个序列。
  \item 本质是债的保全，保障债权的全面实现。\zd{处理逃债}。
  \item \zd{代位请求权}。会结合民诉管辖、诉讼地位、仲裁；经济法税收强制措施考。\al{三角债}。
        \begin{enumerate}
          \item \zd{两货币之债}均合法且到期。\al{甲将自己对丁公司享有的 30\%的股权无偿转让给了丙公司}。
          \item 债务人\zd{怠于行使权利}（或诉或裁）。\al{乙公司多次发函给甲公司要求清偿债务 ×}。
          \item 怠于行权已危害债权人利益。
          \item 不具有人身专属性。
          \item 只能通过诉讼方式行使 \yjt 代位权之诉。
          \item 抗辩权移转。
                \begin{enumerate}
                  \item 债务人对债权人的抗辩。
                  \item 次债务人对债权人的抗辩。次债务人原告，债务人第三人。
                  \item 代位权行使中有的抗辩。\al{行使代位权的数额超出主债权的数额}。
                \end{enumerate}
          \item 物权请求权不能代位，只能请求转让行使。\al{可以代位行使对丙的返还请求权 ×；请求甲将对丙的返还请求权转让给自己 √}。
        \end{enumerate}
  \item \zd{代位保存权/代位权的提前行权}。新增，没考过。
        \begin{enumerate}
          \item 一方未过时效，一方即将过时效/破产。
          \item 未到期，只能入库规则（\zd{提存}）。
        \end{enumerate}
  \item 撤销权。诈害债权。\al{高买低卖、赠与、假离婚}。
        \begin{enumerate}
          \item 要撤的那个节点存在合法有效的债权。
                \begin{enumerate}
                  \item \al{甲 11/6/1 欠乙货款 500 万元，届期无力清偿。10/12/1，甲向丙赠送价值 50W 的机器设备}。
                  \item \zd{不要和商法破产的撤销权搞混了，没说破产用民法的，说了破产用商法的}。
                \end{enumerate}
          \item 主观有逃债恶意，客观上实施了逃债的行为（放弃债权和担保、恶意延长且知道、无偿转让要撤销）。
                \begin{enumerate}
                  \item 恶意串通。\al{乙公司带着礼物去找甲公司}。
                  \item 超过合理对价标准。\al{甲欠乙 100W，房子 90W（市价 120W）卖给丙}。
                  \item 无偿转让。例外。\zd{公证、公益、道德义务赠与不能行使任意撤销权} \yjt 自己不可以撤，债权人可以申请撤。
                        \al{乙公司有权撤销甲公司对戊学校的捐赠}。
                \end{enumerate}
          \item 债务人的行为危害了债权人的利益。
          \item \zd{双重除斥期间}。
        \end{enumerate}
\end{enumerate}
